%% kapa: exam=1 round=23 year=2012 period=2 form=A subject=accounting subject_ko=회계학 seq=18 number=18
%% source: Q-Net 감정평가사 기출문제 (cst00309) artlSeq=5052580 file=raw/1008/5052580/2012년 제23회 감정평가사 제1차시험 기출문제(2교시 A형).pdf
%% conversion: kordoc 4.24.51 -> tools/gen1.py (auto); review=proofread-sonnet; flags=none
\begin{question}{18}
(주)서울은 20×1년 중에 ￦10,100을 지급하고 지분상품을 취득하였는데, 지급액 중 ￦100은 매매수수료이다. 20×1년 말 현재 지분상품의 공정가치는 ￦11,000 이며, (주)서울은 20×2년 초에 지분상품 전체를 ￦11,200에 처분하였다. (주)서울이 이 지분상품을 단기매매금융자산으로 인식할 경우, 이에 대한 회계처리가 20×1년과 20×2년 당기순이익에 미치는 영향은?
\begin{choices}
  \item 20×1년 ￦100 감소, 20×2년 ￦1,200 증가
  \item 20×1년 ￦900 증가, 20×2년 ￦200 증가
  \item 20×1년 ￦900 증가, 20×2년 ￦900 증가
  \item 20×1년 ￦1,000 증가, 20×2년 ￦200 증가
  \item 20×1년 ￦1,000 증가, 20×2년 ￦1,200 증가
\end{choices}
\end{question}
\answer{2}
